\documentclass[11pt]{article}

\usepackage[margin=1.1in]{geometry}
\usepackage{booktabs}
\usepackage{amsmath}
\usepackage[hidelinks]{hyperref}
\usepackage{parskip}

\newcommand{\code}[1]{\texttt{#1}}

\sloppy
\emergencystretch=3em

\title{Timeless dynamic reservation in medios:\\
verification of hold-and-release semantics}
\author{medios research track}
\date{June 9, 2026}

\begin{document}

\maketitle

\begin{center}
\small Companion to \code{docs/research/offers.md}; experiment tracked as
medios issue \#317 (private history),
test landed in \#319 (private history).
\end{center}

\begin{abstract}
medios inherits the full dynamic-reservation machinery from Mesos:
\code{RESERVE} and \code{UNRESERVE} offer operations,
\code{Resource.ReservationInfo}, allocator-side reservation tracking, and
operator-API endpoints. Inherited code in a thin fork is present but not
necessarily functional; large parts of the upstream surface are stubbed, and
the reservation path had no evidence of ever being exercised in medios. We
test the claim from \code{offers.md} that a \emph{timeless} reservation ---
a claim held from now until explicitly released --- is the degenerate case
of the interval-offer model and should work end-to-end on the machinery
already in the tree, with no new product code. A two-stage verification (run
the four inherited reservation tests; add one test for the four properties
they leave uncovered) confirms the claim: resources reserved for a role
persist across allocation cycles with no task running, are never offered to
a framework in another role, are visible in the operator state view, and
return to the shared pool on release. The full suite passes 603/603. The
result validates ``reservation as a claimed resource'' at zero product-code
cost and unblocks the interval lift (lease, expiry, future booking) as a
pure design question.
\end{abstract}

\section{Background and hypothesis}

\code{offers.md} models an offer as a resource vector over a time window,
with the ordinary scalar offer as the degenerate window
$[t_{\mathrm{now}},\, t_{\mathrm{now}} + \mathrm{offer\_timeout}]$. Under
that model a reservation does not add a time dimension; it makes the window
explicit, long, and committed. The minimal reservation is \emph{timeless}:
the window $[\mathrm{now},\, \mathrm{until\ UNRESERVE})$. At that degenerate
point the interval invariant
\begin{equation*}
\sum \mathrm{active\_reservations}(t) \;\le\; \mathrm{capacity}(t)
\end{equation*}
collapses to the static $\mathrm{reserved} + \mathrm{free} =
\mathrm{capacity}$, which is exactly the reserved-versus-unreserved
bucketing the hierarchical allocator already enforces
(\code{RoleTree::trackReservations},
\code{reservationScalarQuantities\_}, \code{hierarchical.cpp:414-458}). If
that reading is correct, medios already ships the timeless reservation and
no allocator work is needed to demonstrate it.

The risk is that the reading is correct about the design but wrong about the
tree. medios strips or stubs much of upstream Mesos (DR-007, DR-013), and a
code path can survive a fork as text without surviving as behavior. The same
session that scoped this experiment found inherited tests that had never run
under medios. Presence had to be separated from function.

\textbf{Hypothesis.} The dynamic-reservation path in the medios tree is
functional end-to-end. Specifically, with no task ever launched:

\begin{itemize}
\item \textbf{H1 (hold).} Resources \code{RESERVE}d for a role persist
  across allocation cycles as allocator state alone --- they do not decay,
  leak back to the free pool, or require a running task to stay claimed.
\item \textbf{H2 (exclusion).} A framework in a different role is never
  offered the reserved resources, even when the cluster is otherwise idle.
\item \textbf{H3 (observability).} The held claim is visible to an operator
  through the controller state view.
\item \textbf{H4 (release).} \code{UNRESERVE} returns the resources to the
  unreserved pool, after which a framework in another role receives them.
\end{itemize}

H1 and H4 together are the hold-and-release contract; H2 is what makes a
reservation a \emph{claim} rather than a preference; H3 is what makes the
claim auditable.

\section{Experiment}

\textbf{Design principle.} Smallest probe first. Stage~0 costs one build and
runs tests that already exist; stage~1 adds a single test and no product
code. A failure at either stage names the broken component directly.

\textbf{Stage 0 --- do the inherited tests pass?} The tree contains four
reservation tests (\code{ReserveCreateLaunchDestroyUnreserve} and
\code{ReserveAndUnreserveResourcesV1}, each parameterized pre/post
reservation-refinement). They were built and run on the Linux build host at
master \code{5b0d21dc3} by the build-dispatch agent.

\textbf{Stage 1 --- close the coverage gap.} Reading the four tests shows
what they do not establish. The framework-path test exercises reserve and
unreserve only from the reserving framework's own point of view; the
operator-API test asserts only HTTP~202 on the calls and never checks that
the reservation took effect. None of H2 or H3 is covered, and H1 is covered
only incidentally. We added one test,
\code{ControllerTest.ReservationCrossRoleHoldAndRelease}
(\code{tests/controller\_tests.cpp}), against an in-process cluster: one
controller, one node (\code{cpus:8;mem:2048}), and two frameworks --- A in
role \code{alpha}, B in role \code{beta}. The choreography:

\begin{enumerate}
\item A receives the full node and applies \code{RESERVE} \code{cpus:4} for
  \code{alpha} (dynamic reservation, principal matching A's).
\item A declines its offers for three consecutive allocation cycles,
  asserting each time that the re-offer still contains the reserved
  \code{cpus:4} (H1).
\item The controller \code{/state} endpoint is asserted to report
  \code{nodes[0].reserved\_resources.alpha.cpus == 4} (H3).
\item A suppresses; B starts and receives offers. B's offer is asserted to
  contain only the unreserved remainder (\code{cpus == 4}) and no resource
  reserved to \code{alpha} (H2).
\item B suppresses; A revives, receives the still-held reservation, applies
  \code{UNRESERVE}, and suppresses again.
\item B revives and is asserted to receive the full node, all unreserved
  (\code{cpus:8;mem:2048} allocated to \code{beta}); \code{/state} is
  asserted to show the \code{alpha} reservation gone and
  \code{unreserved\_resources.cpus == 8} (H4).
\end{enumerate}

Only one framework is eligible for offers at each step (suppress/revive
choreography with queue-drain barriers between phases), so every assertion
is against a deterministic allocation, not a race. The probe uses
\code{cpus} rather than the issue's \code{gpu:4}: GPU resources are gated
behind a separate framework capability (\code{GPU\_RESOURCES},
\code{hierarchical.cpp:3113}) that is orthogonal to reservation semantics,
and any scalar exercises the same reserve/track/offer/release path.

\section{Results}

All stages pass on the Linux build host (full \code{ctest} suite,
build-dispatch agent):

\begin{center}
\begin{tabular}{lll}
\toprule
Stage & Test & Result \\
\midrule
0 & \code{ReserveCreateLaunchDestroyUnreserve} (pre \& post) & pass, 0.19s each \\
0 & \code{ReserveAndUnreserveResourcesV1} (pre \& post) & pass, 0.15s each \\
1 & \code{ReservationCrossRoleHoldAndRelease} (H1--H4) & pass, 0.30s \\
--- & full suite at \code{ac3ccc7ac} & 603/603 \\
\bottomrule
\end{tabular}
\end{center}

Secondary findings from the verification:

\begin{itemize}
\item \textbf{Observability is per-node only.} \code{/state} and
  \code{GET\_AGENTS} expose reservations on each node, which satisfies H3.
  The \code{/roles} endpoint does not report reserved-but-unallocated
  quantities per role, and the allocator's aggregate
  (\code{reservationScalarQuantities\_}) is not surfaced over HTTP. An
  operator can see \emph{where} resources are held but has no single
  cluster-wide reserved-versus-free figure per role.
\item \textbf{The inherited operator-API test is status-code-only.} It would
  pass even if \code{RESERVE} were a silent no-op. Stage~1 closes that hole
  through the framework path; the operator path's effect is still asserted
  only indirectly.
\end{itemize}

\textbf{Threats to validity.} One node, one resource shape, scalar resources
only. The test does not exercise node restart (reservation checkpoint
recovery), controller failover, reservation refinement under hierarchical
roles, principal/ACL enforcement matrices, or concurrent reservation
conflicts. It verifies the contract, not its hardening.

\section{Conclusion}

The hypothesis holds on all four properties. The timeless dynamic
reservation --- claim held until released, no time dimension --- works
end-to-end on machinery already in the medios tree, at zero product-code
cost. This validates the \code{offers.md} framing empirically: the
reserved/unreserved bucketing is the degenerate interval invariant, and
medios ships the degenerate reservation today. The fork's ``present but
possibly stubbed'' risk was real in general but did not materialize on this
path.

\section{Next steps}

\begin{enumerate}
\item \textbf{Interval lift (v2).} Lease/expiry (auto-\code{UNRESERVE} at
  $t_1$) and future booking (reserve from future $t_0$) are now a design
  exercise on a verified base. Admission against an uncertain future and
  the reserved-but-could-not-be-honored settlement state (\code{offers.md}
  open questions) are the substantive design work.
\item \textbf{Gang feasibility
  (\#316, private history).} The
  placement simulation is independent of this result but shares the
  research goal: keep the allocator greedy. Its v2 stage composes with
  reservations verified here.
\item \textbf{Per-role reserved aggregate.} Small observability lift:
  surface reserved-versus-free per role (the allocator already tracks it).
  Worth doing when an operator-facing consumer exists; not before.
\item \textbf{Operator-path effect assertion.} Extend the V1 operator-API
  test to assert the reservation's effect, not just the 202.
\item \textbf{SDK surface.} Reservation wire types and a worked example in
  the Go SDK stay out of scope until the protocol shape (including v2) is
  settled.
\end{enumerate}

\end{document}
